% ====================================================================
% ФАЙЛ: tasks/01_mechanics/gravity_and_newton3_bank.tex
% ТЕМА: ЗАКОН ВСЕМИРНОГО ТЯГОТЕНИЯ И ТРЕТИЙ ЗАКОН НЬЮТОНА
% ПАЧКА 1: ВАРИАНТЫ 1–10
% ====================================================================

% === ЗАДАЧА 1 ===
% Тег: #МЕХАНИКА #КИМ_06 #ВАРИАНТ_01
\begin{taskbasic}[code={\label{task:gravity_v1_n6}}]
\textbf{Задача 1.} Спутник обращается вокруг Земли по круговой орбите. В результате перехода на другую круговую орбиту центростремительное ускорение спутника уменьшается. Как изменяются при этом переходе радиус орбиты спутника и его скорость движения по орбите?

Для каждой величины определите соответствующий характер изменения:
1) увеличивается;
2) уменьшается;
3) не изменяется.

Запишите в таблицу выбранные цифры для каждой физической величины. Цифры в ответе могут повторяться.
\begin{center}
\begin{tabular}{|c|c|}
\hline
Радиус орбиты & Скорость спутника \\ \hline
 & \\ \hline
\end{tabular}
\end{center}

\kim{6}
\end{taskbasic}
\answer{12}

% === ЗАДАЧА 2 ===
% Тег: #МЕХАНИКА #КИМ_06 #ВАРИАНТ_02
\begin{taskbasic}[code={\label{task:gravity_v2_n6}}]
\textbf{Задача 2.} Спутник обращается вокруг Земли по круговой орбите. В результате перехода на другую круговую орбиту скорость движения спутника увеличивается. Как изменяются при этом переходе радиус орбиты спутника и его период обращения вокруг Земли?

Для каждой величины определите соответствующий характер изменения:
1) увеличивается;
2) уменьшается;
3) не изменяется.

Запишите в таблицу выбранные цифры для каждой физической величины. Цифры в ответе могут повторяться.
\begin{center}
\begin{tabular}{|c|c|}
\hline
Радиус орбиты & Период обращения \\ \hline
 & \\ \hline
\end{tabular}
\end{center}

\kim{6}
\end{taskbasic}
\answer{22}

% === ЗАДАЧА 3 ===
% Тег: #МЕХАНИКА #КИМ_02 #ВАРИАНТ_03
\begin{taskbasic}[code={\label{task:gravity_v3_n2}}]
\textbf{Задача 3.} Две планеты с одинаковыми массами вращаются по круговым орбитам вокруг звезды. Петя установил, что сила притяжения первой планеты к звезде в $4$ раза больше, чем сила притяжения второй планеты. Чему равно отношение радиусов орбит этих планет $\frac{R_2}{R_1}$?

\kim{2}
\end{taskbasic}
\answer{2}

% === ЗАДАЧА 4 ===
% Тег: #МЕХАНИКА #КИМ_02 #ВАРИАНТ_04
\begin{taskbasic}[code={\label{task:gravity_v4_n2}}]
\textbf{Задача 4.} Две планеты вращаются по круговым орбитам вокруг звезды. Радиус орбиты первой планеты в $2$ раза больше радиуса орбиты второй планеты, а масса первой планеты в $2$ раза меньше массы второй планеты. Чему равно отношение сил притяжения планет к звезде $\frac{F_1}{F_2}$?

\kim{2}
\end{taskbasic}
\answer{0{,}125}

% === ЗАДАЧА 5 ===
% Тег: #МЕХАНИКА #КИМ_06 #ВАРИАНТ_05
\begin{taskbasic}[code={\label{task:gravity_v5_n6}}]
\textbf{Задача 5.} Искусственный спутник обращается вокруг Земли по круговой орбите. Радиус этой орбиты уменьшили. Как изменились при этом переходе модуль силы тяжести, действующей на спутник, и его кинетическая энергия?

Для каждой величины определите соответствующий характер изменения:
1) увеличилась;
2) уменьшилась;
3) не изменилась.

Запишите в таблицу выбранные цифры для каждой физической величины. Цифры в ответе могут повторяться.
\begin{center}
\begin{tabular}{|c|c|}
\hline
Сила тяжести & Кинетическая энергия \\ \hline
 & \\ \hline
\end{tabular}
\end{center}

\kim{6}
\end{taskbasic}
\answer{11}

% === ЗАДАЧА 6 ===
% Тег: #МЕХАНИКА #КИМ_06 #ВАРИАНТ_06
\begin{taskbasic}[code={\label{task:gravity_v6_n6}}]
\textbf{Задача 6.} Искусственный спутник обращается вокруг Земли по круговой орбите. Радиус этой орбиты увеличили. Как изменились при этом переходе период обращения спутника и его потенциальная энергия в системе «Спутник — Земля»?

Для каждой величины определите соответствующий характер изменения:
1) увеличилась;
2) уменьшилась;
3) не изменилась.

Запишите в таблицу выбранные цифры для каждой физической величины. Цифры в ответе могут повторяться.
\begin{center}
\begin{tabular}{|c|c|}
\hline
Период обращения & Потенциальная энергия \\ \hline
 & \\ \hline
\end{tabular}
\end{center}

\kim{6}
\end{taskbasic}
\answer{11}

% === ЗАДАЧА 7 ===
% Тег: #МЕХАНИКА #КИМ_02 #ВАРИАНТ_07
\begin{taskbasic}[code={\label{task:gravity_v7_n2}}]
\textbf{Задача 7.} Две материальные точки находятся на расстоянии $R$ друг от друга и притягиваются друг к другу с силой $F$. Во сколько раз увеличится сила гравитационного притяжения между ними, если массу одной из точек увеличить в $5$ раз, а расстояние между ними увеличить в $2$ раза?

\kim{2}
\end{taskbasic}
\answer{1{,}25}

% === ЗАДАЧА 8 ===
% Тег: #МЕХАНИКА #КИМ_02 #ВАРИАНТ_08
\begin{taskbasic}[code={\label{task:gravity_v8_n2}}]
\textbf{Задача 8.} Две материальные точки массами $m_1$ и $m_2$ находятся на расстоянии $R$ друг от друга и притягиваются гравитационными силами. Во сколько раз уменьшится сила гравитационного притяжения между ними, если расстояние между ними увеличить в $4$ раза, а массу одной из точек увеличить в $4$ раза?

\kim{2}
\end{taskbasic}
\answer{4}

% === ЗАДАЧА 9 ===
% Тег: #МЕХАНИКА #КИМ_06 #ВАРИАНТ_09
\begin{taskbasic}[code={\label{task:gravity_v9_n6}}]
\textbf{Задача 9.} Планета движется вокруг звезды по круговой орбите. В результате перехода на другую круговую орбиту скорость её движения уменьшается. Как изменяются при этом переходе период обращения планеты и её центростремительное ускорение?

Для каждой величины определите соответствующий характер изменения:
1) увеличивается;
2) уменьшается;
3) не изменяется.

Запишите в таблицу выбранные цифры для каждой физической величины. Цифры в ответе могут повторяться.
\begin{center}
\begin{tabular}{|c|c|}
\hline
Период обращения & Центростремительное ускорение \\ \hline
 & \\ \hline
\end{tabular}
\end{center}

\kim{6}
\end{taskbasic}
\answer{12}

% === ЗАДАЧА 10 ===
% Тег: #МЕХАНИКА #КИМ_06 #ВАРИАНТ_10
\begin{taskbasic}[code={\label{task:gravity_v10_n6}}]
\textbf{Задача 10.} Планета движется вокруг звезды по круговой орбите. В результате перехода на другую круговую орбиту период её обращения уменьшается. Как изменяются при этом переходе радиус орбиты планеты и её скорость движения по орбите?

Для каждой величины определите соответствующий характер изменения:
1) увеличивается;
2) уменьшается;
3) не изменяется.

Запишите в таблицу выбранные цифры для каждой физической величины. Цифры в ответе могут повторяться.
\begin{center}
\begin{tabular}{|c|c|}
\hline
Радиус орбиты & Скорость планеты \\ \hline
 & \\ \hline
\end{tabular}
\end{center}

\kim{6}
\end{taskbasic}
\answer{21}

% ====================================================================
% ФАЙЛ: tasks/01_mechanics/gravity_and_newton3_bank_part2.tex
% ТЕМА: ЗАКОН ВСЕМИРНОГО ТЯГОТЕНИЯ И ТРЕТИЙ ЗАКОН НЬЮТОНА
% ПАЧКА 2: ВАРИАНТЫ 11–20
% ====================================================================

% === ЗАДАЧА 1 ===
% Тег: #МЕХАНИКА #КИМ_06 #ВАРИАНТ_13
\begin{taskbasic}[code={\label{task:gravity_v13_n6}}]
\textbf{Задача 1.} Космический аппарат, обращающийся вокруг Венеры по круговой орбите, перешёл на другую круговую орбиту большего радиуса. Как изменились в результате этого перехода центростремительное ускорение, с которым аппарат движется по орбите, и его период обращения вокруг Венеры?

Для каждой величины определите соответствующий характер изменения:
1) увеличилось;
2) уменьшилось;
3) не изменилось.

Запишите в таблицу выбранные цифры для каждой физической величины. Цифры в ответе могут повторяться.
\begin{center}
\begin{tabular}{|c|c|}
\hline
Центростремительное ускорение & Период обращения \\ \hline
 & \\ \hline
\end{tabular}
\end{center}

\kim{6}
\end{taskbasic}
\answer{21}

% === ЗАДАЧА 2 ===
% Тег: #МЕХАНИКА #КИМ_06 #ВАРИАНТ_14
\begin{taskbasic}[code={\label{task:gravity_v14_n6}}]
\textbf{Задача 2.} Космический аппарат, обращающийся вокруг Луны по круговой орбите, перешёл на другую круговую орбиту меньшего радиуса. Как изменились в результате этого перехода кинетическая энергия аппарата и его частота обращения вокруг Луны?

Для каждой величины определите соответствующий характер изменения:
1) увеличилась;
2) уменьшилась;
3) не изменилась.

Запишите в таблицу выбранные цифры для каждой физической величины. Цифры в ответе могут повторяться.
\begin{center}
\begin{tabular}{|c|c|}
\hline
Кинетическая энергия & Частота обращения \\ \hline
 & \\ \hline
\end{tabular}
\end{center}

\kim{6}
\end{taskbasic}
\answer{11}

% === ЗАДАЧА 3 ===
% Тег: #МЕХАНИКА #КИМ_02 #ВАРИАНТ_15
\begin{taskbasic}[code={\label{task:gravity_v15_n2}}]
\textbf{Задача 3.} Две материальные точки находятся на расстоянии $R$ друг от друга и притягиваются друг к другу гравитационными силами с модулем $F$. Во сколько раз уменьшится модуль сил гравитационного притяжения между ними, если расстояние между ними увеличить в $3$ раза, а массу одной из точек увеличить в $3$ раза?

\kim{2}
\end{taskbasic}
\answer{3}

% === ЗАДАЧА 4 ===
% Тег: #МЕХАНИКА #КИМ_02 #ВАРИАНТ_16
\begin{taskbasic}[code={\label{task:gravity_v16_n2}}]
\textbf{Задача 4.} Две материальные точки находятся на расстоянии $R$ друг от друга и притягиваются друг к другу гравитационными силами с модулем $F$. Во сколько раз увеличится модуль сил гравитационного притяжения между ними, если массу одной из точек увеличить в $6$ раз, а расстояние между ними увеличить в $2$ раза?

\kim{2}
\end{taskbasic}
\answer{1{,}5}

% === ЗАДАЧА 5 ===
% Тег: #МЕХАНИКА #КИМ_05 #ВАРИАНТ_17
\begin{taskbasic}[code={\label{task:gravity_v17_n5}}]
\textbf{Задача 5.} Используя таблицу, содержащую сведения о планетах Солнечной системы, выберите все верные утверждения, соответствующие характеристикам планет.
\begin{center}
\begin{tabular}{|l|c|c|c|}
\hline
\textbf{Планета} & \textbf{Средний радиус, км} & \textbf{Масса, $\cdot 10^{24}$~кг} & \textbf{Период вращения вокруг оси} \\ \hline
Меркурий & $2440$ & $0{,}33$ & $58{,}6$~сут. \\ \hline
Венера   & $6052$ & $4{,}87$ & $243$~сут. \\ \hline
Земля    & $6371$ & $5{,}97$ & $23$~ч $56$~мин $4$~с \\ \hline
Марс     & $3390$ & $0{,}642$ & $24$~ч $37$~мин $23$~с \\ \hline
\end{tabular}
\end{center}

1) Ускорение свободного падения на Венере составляет примерно $8{,}9$~м/с$^{2}$.
2) Масса Марса примерно в $10$ раза больше массы Земли.
3) Меркурий делает один оборот вокруг своей оси быстрее, чем Земля.
4) Сила гравитационного притяжения Венеры к Солнцу больше, чем у Земли.
5) Первая космическая скорость для Марса составляет примерно $3{,}6$~км/ч.

\kim{5}
\end{taskbasic}
\answer{15}

% === ЗАДАЧА 6 ===
% Тег: #МЕХАНИКА #КИМ_05 #ВАРИАНТ_18
\begin{taskbasic}[code={\label{task:gravity_v18_n5}}]
\textbf{Задача 6.} Используя таблицу, содержащую сведения о планетах Солнечной системы, выберите все верные утверждения, соответствующие характеристикам планет.
\begin{center}
\begin{tabular}{|l|c|c|c|}
\hline
\textbf{Планета} & \textbf{Средний радиус, км} & \textbf{Масса, $\cdot 10^{24}$~кг} & \textbf{Период вращения вокруг оси} \\ \hline
Меркурий & $2440$ & $0{,}33$ & $58{,}6$~сут. \\ \hline
Венера   & $6052$ & $4{,}87$ & $243$~сут. \\ \hline
Земля    & $6371$ & $5{,}97$ & $23$~ч $56$~мин $4$~с \\ \hline
Марс     & $3390$ & $0{,}642$ & $24$~ч $37$~мин $23$~с \\ \hline
\end{tabular}
\end{center}

1) Ускорение свободного падения на Меркурии составляет примерно $5{,}8$~м/с$^{2}$.
2) Венера делает один оборот вокруг своей оси быстрее, чем Марс.
3) Первая космическая скорость для Земли составляет примерно $7{,}9$~км/с.
4) Ускорение свободного падения на Марсе меньше, чем на Земле.
5) Первая космическая скорость для Меркурия больше, чем для Венеры.

\kim{5}
\end{taskbasic}
\answer{34}

% ====================================================================
% ФАЙЛ: tasks/01_mechanics/gravity_and_newton3_bank_part3.tex
% ТЕМА: ЗАКОН ВСЕМИРНОГО ТЯГОТЕНИЯ И ТРЕТИЙ ЗАКОН НЬЮТОНА
% ПАЧКА 3: ВАРИАНТЫ 21–30
% ====================================================================

% === ЗАДАЧА 1 ===
% Тег: #МЕХАНИКА #КИМ_02 #ВАРИАНТ_23
\begin{taskbasic}[code={\label{task:gravity_v23_n2}}]
\textbf{Задача 1.} Модуль сил гравитационного притяжения между двумя шарами, находящимися на расстоянии $R$ друг от друга, равен $16$~нН. Каков будет модуль сил притяжения между ними, если расстояние между ними увеличить в $4$ раза?

\kim{2}
\end{taskbasic}
\answer{1~нН}

% === ЗАДАЧА 2 ===
% Тег: #МЕХАНИКА #КИМ_02 #ВАРИАНТ_24
\begin{taskbasic}[code={\label{task:gravity_v24_n2}}]
\textbf{Задача 2.} Модуль сил гравитационного притяжения между двумя шарами, находящимися на расстоянии $2$~м друг от друга, равен $9$~нН. Каков будет модуль сил притяжения между ними, если расстояние увеличить до $6$~м?

\kim{2}
\end{taskbasic}
\answer{1~нН}

% === ЗАДАЧА 3 ===
% Тег: #МЕХАНИКА #КИМ_02 #ВАРИАНТ_25
\begin{taskbasic}[code={\label{task:gravity_v25_n2}}]
\textbf{Задача 3.} На графике приведена зависимость ускорения $a$ бруска, скользящего без трения по горизонтальной поверхности, от величины приложенной к нему горизонтальной силы $\vec{F}$. Систему отсчёта считать инерциальной. Чему равна масса бруска?
\begin{center}
\begin{tikzpicture}[x=0.6cm, y=10cm, every node/.style={font=\footnotesize}]
    \draw[xstep=1, ystep=0.05, gray!30, thin] (0,0) grid (6,0.25);
    \draw[-{Stealth[scale=0.8]}, black, thick] (0,0) -- (6.5,0) node[right] {$F, \text{ Н}$};
    \draw[-{Stealth[scale=0.8]}, black, thick] (0,0) -- (0,0.27) node[above] {$a, \text{ м/с}^{2}$};
    \foreach \x in {2,4,6} \node[below] at (\x,0) {\x};
    \node[left] at (0,0.1) {$0{,}1$};
    \node[left] at (0,0.2) {$0{,}2$};
    \node[below left] at (0,0) {0};
    \draw[ultra thick, brandPrimary] (0,0) -- (5,0.2);
    \fill[black] (5,0.2) circle (1.5pt);
\end{tikzpicture}
\end{center}

\kim{2}
\end{taskbasic}
\answer{25~кг}

% === ЗАДАЧА 4 ===
% Тег: #МЕХАНИКА #КИМ_02 #ВАРИАНТ_26
\begin{taskbasic}[code={\label{task:gravity_v26_n2}}]
\textbf{Задача 4.} На графике приведена зависимость модуля горизонтальной силы $\vec{F}$, приложенной к бруску, скользящего без трения по горизонтальной поверхности, от модуля его ускорения $a$. Систему отсчёта считать инерциальной. Чему равна масса бруска?
\begin{center}
\begin{tikzpicture}[x=6cm, y=0.25cm, every node/.style={font=\footnotesize}]
    \draw[xstep=0.1, ystep=2, gray!30, thin] (0,0) grid (0.6,12);
    \draw[-{Stealth[scale=0.8]}, black, thick] (0,0) -- (0.68,0) node[right] {$a, \text{ м/с}^{2}$};
    \draw[-{Stealth[scale=0.8]}, black, thick] (0,0) -- (0,13) node[above] {$F, \text{ Н}$};
    \foreach \y in {3,6,12} \node[left] at (0,\y) {\y};
    \node[below] at (0.2,0) {$0{,}2$};
    \node[below] at (0.4,0) {$0{,}4$};
    \node[below] at (0.6,0) {$0{,}6$};
    \node[below left] at (0,0) {0};
    \draw[ultra thick, brandPrimary] (0,0) -- (0.6,12);
    \fill[black] (0.6,12) circle (1.5pt);
\end{tikzpicture}
\end{center}

\kim{2}
\end{taskbasic}
\answer{20~кг}